\documentclass[11pt,a4paper]{article}

\usepackage{amsmath,amssymb,amsthm}
\usepackage{algorithm}
\usepackage{algpseudocode}
\usepackage{listings}
\usepackage{hyperref}
\usepackage{booktabs}
\usepackage{geometry}
\usepackage{graphicx}
\usepackage{xcolor}

\geometry{margin=2.5cm}

\lstset{
  language=C++,
  basicstyle=\ttfamily\small,
  keywordstyle=\bfseries\color{blue!70!black},
  commentstyle=\itshape\color{green!50!black},
  numbers=left,
  numberstyle=\tiny\color{gray},
  frame=single,
  breaklines=true,
  columns=flexible
}

\newtheorem{definition}{Definition}
\newtheorem{lemma}{Lemma}
\newtheorem{theorem}{Theorem}
\newtheorem{remark}{Remark}

\title{%
  Triple-Double Floating-Point Arithmetic: \\
  Algorithms and Implementation in the QD Library%
}
\author{%
  Nakata Maho \\
  \textit{RIKEN} \\
  \texttt{maho@riken.jp}%
}
\date{2026}

\begin{document}
\maketitle

\begin{abstract}
We describe the design and algorithmic details of \texttt{td\_real}, a
triple-double floating-point type that extends the QD library
(Hida--Li--Bailey) with approximately 159-bit ($\approx 47$ decimal
digit) precision.  A triple-double number is represented as an
unevaluated sum of three IEEE 754 double-precision values
$a = x_0 + x_1 + x_2$, subject to the non-overlapping constraint
$|x_{i+1}| \le \tfrac{1}{2}\,\mathrm{ulp}(x_i)$.  We present
algorithms for renormalization, addition (both IEEE-style merge and
Knuth--M{\o}ller chain), subtraction, multiplication, division with
overflow-safe rescaling, square root via Newton iteration, and
transcendental functions (exp, log, sin, cos, atan2, sinh, cosh,
etc.).  We also describe the type-conversion algorithms and
mixed-precision operator semantics that enable seamless
interoperability among DD, TD, and QD types, including promotion
rules, lossy truncation with renormalization, and compound-assignment
handling.  A detailed binary64 flop count analysis shows that TD
achieves $\approx$74\% of QD's precision at $\approx$49\% of
QD's addition cost and $\approx$64\% of its multiplication cost.
C and Fortran 90 binding layers are also described.  The
implementation is publicly available under a BSD 2-clause license.
\end{abstract}

\section{Introduction}

Many applications in computational science require more than the
$\approx 16$ decimal digits offered by IEEE 754 binary64
(double-precision).  The double-double (DD, $\approx 106$-bit,
$\approx 31$-digit) and quad-double (QD, $\approx 212$-bit,
$\approx 63$-digit) types introduced by Hida, Li, and
Bailey~\cite{HLB2001,QDlib} fill a useful niche between hardware
double and arbitrary-precision libraries such as GMP/MPFR.  However,
the factor-of-two jump in both precision and cost from DD to QD
leaves a gap: problems that need 40--50 digits pay the full cost of
QD's 63 digits.

The \texttt{td\_real} type fills this gap.  It stores three
non-overlapping doubles, yielding $\approx 159$ bits ($\approx 47$
decimal digits) of significand.  As shown in Section~\ref{sec:cost},
TD addition (default sloppy variant) costs 41~flops vs.\ QD's~84
(49\%), and TD multiplication costs 116/206~flops (FMA/no FMA)
vs.\ QD's 182/332 (64\%/62\%).

This paper is organized as follows.  Section~\ref{sec:eft}
introduces the error-free transformation (EFT) primitives shared
with the DD and QD libraries.
Section~\ref{sec:representation}
defines the triple-double representation and its invariants.
Section~\ref{sec:renorm} presents the renormalization algorithms.
Sections~\ref{sec:add}--\ref{sec:div} describe the four basic
arithmetic operations.  Section~\ref{sec:cost} gives a detailed
binary64 flop count comparison among DD, TD, and QD.
Section~\ref{sec:mixed} details the
type-conversion algorithms and mixed-precision operator semantics
for interoperability among DD, TD, and QD.
Section~\ref{sec:sqrt} covers square root.
Section~\ref{sec:transcendental} describes the transcendental
function implementations.  Section~\ref{sec:constants} lists the
precomputed constants.  Section~\ref{sec:bindings} outlines the C
and Fortran binding layers.
Section~\ref{sec:conclusion} concludes.


\section{Preliminaries: Error-Free Transformations}
\label{sec:eft}

The triple-double arithmetic, like the DD and QD types before it,
is built upon a small set of \emph{error-free transformations}
(EFTs)---operations that compute both the floating-point result and
its rounding error exactly in double precision.  These primitives
reside in \texttt{include/inline.h} within the \texttt{qd} namespace
and are shared by all three multi-word types.

\begin{definition}[\textsc{TwoSum}~{\cite{Knuth1997,Dekker1971}}]
Given IEEE 754 doubles $a$ and $b$, \textsc{TwoSum} computes
$s = \mathrm{fl}(a + b)$ and $e = (a + b) - s$ exactly, so that
$a + b = s + e$:
\begin{lstlisting}
double two_sum(double a, double b, double &err) {
  double s = a + b;
  double bb = s - a;
  err = (a - (s - bb)) + (b - bb);
  return s;
}
\end{lstlisting}
Cost: 6 flops.  No assumption on the relative magnitudes of $a$
and $b$.
\end{definition}

\begin{definition}[\textsc{QuickTwoSum}~{\cite{Dekker1971}}]
A faster variant requiring $|a| \ge |b|$:
\begin{lstlisting}
double quick_two_sum(double a, double b, double &err) {
  double s = a + b;
  err = b - (s - a);
  return s;
}
\end{lstlisting}
Cost: 3 flops.  Used extensively in renormalization where the
magnitude ordering is known.
\end{definition}

\begin{definition}[\textsc{TwoDiff}]
The subtraction analogue of \textsc{TwoSum}: computes
$s = \mathrm{fl}(a - b)$ and $e$ such that $a - b = s + e$.
\begin{lstlisting}
double two_diff(double a, double b, double &err) {
  double s = a - b;
  double bb = s - a;
  err = (a - (s - bb)) - (b + bb);
  return s;
}
\end{lstlisting}
\end{definition}

\begin{definition}[\textsc{TwoProd}~{\cite{Dekker1971}}]
Computes $p = \mathrm{fl}(a \cdot b)$ and $e = a \cdot b - p$
exactly.  When a fused multiply-add (FMA) instruction is available
(\texttt{QD\_FMS} defined), this reduces to two operations:
\begin{lstlisting}
double two_prod(double a, double b, double &err) {
#ifdef QD_FMS
  double p = a * b;
  err = QD_FMS(a, b, p);   // fma(a, b, -p)
  return p;
#else
  // Dekker splitting fallback: 17 flops
  double a_hi, a_lo, b_hi, b_lo;
  double p = a * b;
  split(a, a_hi, a_lo);
  split(b, b_hi, b_lo);
  err = ((a_hi*b_hi - p) + a_hi*b_lo + a_lo*b_hi) + a_lo*b_lo;
  return p;
#endif
}
\end{lstlisting}
The \texttt{split} function uses the Veltkamp splitting constant
$\sigma = 2^{27} + 1$ to decompose a double into 26-bit halves,
with a threshold guard at $2^{996}$ to avoid overflow.
\end{definition}

\begin{definition}[\textsc{ThreeSum}]
A TD-specific helper that exactly adds three doubles $a, b, c$ and
produces three non-overlapping outputs:
\begin{lstlisting}
void three_sum(double &a, double &b, double &c) {
  double t1, t2, t3;
  t1 = two_sum(a, b, t2);
  a  = two_sum(c, t1, t3);
  b  = two_sum(t2, t3, c);
}
\end{lstlisting}
Cost: 3 \textsc{TwoSum} = 18 flops.  After the call,
$a_{\text{old}} + b_{\text{old}} + c_{\text{old}} = a + b + c$
with $|b| \le \frac{1}{2}\,\mathrm{ulp}(a)$ and
$|c| \le \frac{1}{2}\,\mathrm{ulp}(b)$.  This is used in the
TD$\times$TD multiplication to accumulate partial products.
\end{definition}

\begin{definition}[\textsc{QuickThreeAccum}]
Absorbs a new value $c$ into a running pair $(a, b)$ and returns
a ``finished'' high word when both $a$ and $b$ are non-zero:
\begin{lstlisting}
double quick_three_accum(double &a, double &b, double c) {
  double s = two_sum(b, c, b);
  s = two_sum(a, s, a);
  if (a != 0.0 && b != 0.0) return s;
  if (b == 0.0) { b = a; a = s; }
  else           { a = s; }
  return 0.0;
}
\end{lstlisting}
This is the core of the IEEE-style merge addition
(Section~\ref{sec:add}).  Returning 0 signals that the accumulator
has absorbed the value without emitting a resolved word.
\end{definition}

\begin{definition}[\textsc{MulPwr2}]
Multiplies a triple-double by a power of two, component-wise:
\begin{lstlisting}
td_real mul_pwr2(const td_real &a, double d) {
  return td_real(a[0]*d, a[1]*d, a[2]*d);
}
\end{lstlisting}
Since $d$ is a power of two, each multiplication is exact (no
rounding).  This operation is $O(1)$ and is used extensively in
the overflow-safe rescaling of division and square root
(Sections~\ref{sec:div} and~\ref{sec:sqrt}).
\end{definition}

\begin{remark}[The \textsc{Split} threshold]
\label{rem:split}
The \textsc{TwoProd} non-FMA path uses the Veltkamp splitting
constant $\sigma = 2^{27} + 1$.  Computing $\sigma \cdot a$ can
overflow if $|a| > 2^{996}$ (since
$2^{996} \cdot (2^{27} + 1) \approx 2^{1023}$).  The
\texttt{split} function therefore prescales by $2^{-28}$ when
$|a| > 2^{996}$ and post-scales by $2^{28}$.  This built-in
threshold protects \emph{all} EFT-based arithmetic, but division
and square root require additional rescaling because their
intermediate computations can generate values that challenge the
$2^{996}$ boundary---see Remark~\ref{rem:threshold969}.
\end{remark}

These seven primitives---\textsc{TwoSum}, \textsc{QuickTwoSum},
\textsc{TwoDiff}, \textsc{TwoProd}, \textsc{ThreeSum},
\textsc{QuickThreeAccum}, and \textsc{MulPwr2}---together with
the \textsc{Split} threshold mechanism, form the complete
foundation upon which all TD arithmetic operations are built.

Table~\ref{tab:eft-costs} summarizes the cost of each primitive
in binary64 floating-point operations.  Since only \textsc{TwoProd}
involves a multiplication with error tracking, the FMA
availability affects it alone; all other primitives are
addition-based and have identical costs in both cases.

\begin{table}[h]
\centering
\caption{Cost of EFT primitives in binary64 flops.
  FMA = 1 flop.}
\label{tab:eft-costs}
\begin{tabular}{lccl}
\toprule
Primitive & FMA & No FMA & Notes \\
\midrule
\textsc{TwoSum}          & 6  & 6  & Knuth/Dekker \\
\textsc{QuickTwoSum}     & 3  & 3  & requires $|a|\ge|b|$ \\
\textsc{TwoDiff}         & 6  & 6  & subtraction analogue \\
\textsc{TwoProd}         & 2  & 17 & FMA: $p$+fms; else: Dekker split \\
\textsc{ThreeSum}        & 18 & 18 & $= 3\times\textsc{TwoSum}$ \\
\textsc{ThreeSum2}       & 13 & 13 & $= 2\times\textsc{TwoSum}+1$ \\
\textsc{QuickThreeAccum} & 12 & 12 & $= 2\times\textsc{TwoSum}$ \\
\textsc{MulPwr2}         & 3  & 3  & component-wise, exact \\
\bottomrule
\end{tabular}
\end{table}


\section{Representation and Invariants}
\label{sec:representation}

\begin{definition}[Triple-double number]
A triple-double number $a$ is an unevaluated sum of three IEEE 754
binary64 values:
\begin{equation}
  a = x_0 + x_1 + x_2, \quad |x_1| \le \tfrac{1}{2}\,\mathrm{ulp}(x_0),
  \quad |x_2| \le \tfrac{1}{2}\,\mathrm{ulp}(x_1).
  \label{eq:td-rep}
\end{equation}
The component $x_0$ carries the leading bits; $x_1$ and $x_2$ carry
successively lower-order bits.  The struct layout is:
\begin{lstlisting}
struct td_real {
  double x[3];  // x[0]: high, x[1]: mid, x[2]: low
};
\end{lstlisting}
\end{definition}

\begin{remark}[Precision]
Each double carries 53 bits of significand.  In the ideal
non-overlapping case, the triple-double significand spans
$53 + 52 + 52 = 157$ bits, giving
$\lfloor 157 \log_{10} 2 \rfloor = 47$ decimal digits.
The machine epsilon is $\varepsilon_{\mathrm{td}} = 2^{-157}
\approx 5.474 \times 10^{-48}$.
\end{remark}

\begin{remark}[Relationship to DD and QD]
The DD type uses 2 doubles ($\approx 106$ bits); QD uses 4 doubles
($\approx 212$ bits).  TD sits between them:
\begin{center}
\begin{tabular}{lccc}
\toprule
Type & Components & Bits & Decimal digits \\
\midrule
\texttt{dd\_real} & 2 & $\approx 106$ & 31 \\
\texttt{td\_real} & 3 & $\approx 157$ & 47 \\
\texttt{qd\_real} & 4 & $\approx 212$ & 63 \\
\bottomrule
\end{tabular}
\end{center}
\end{remark}

\subsection{Special values}

The implementation represents $\pm\infty$ by setting all three
components to $\pm\infty$, and NaN by setting all three components to
NaN.  Membership tests use the macros \texttt{QD\_ISINF} and
\texttt{QD\_ISNAN} on the leading component $x_0$, which suffices
because the non-overlapping invariant guarantees $|x_1| \ll |x_0|$
for finite values:
\begin{lstlisting}
bool isnan()    const { return QD_ISNAN(x[0]) || QD_ISNAN(x[1]) || QD_ISNAN(x[2]); }
bool isfinite() const { return QD_ISFINITE(x[0]); }
bool isinf()    const { return QD_ISINF(x[0]); }
\end{lstlisting}


\section{Renormalization}
\label{sec:renorm}

After each arithmetic operation the result typically consists of more
than three double-precision components that may overlap.  The
\emph{renormalization} step merges and sorts these components so that
the non-overlapping invariant~\eqref{eq:td-rep} is restored.

Three overloads are provided: for 3, 4, and 5 input components.

\subsection{\texttt{renorm(c0, c1, c2)} --- 3-component}

This is the simplest case.  The algorithm uses \textsc{QuickTwoSum}
to propagate carries from high to low:
\begin{algorithmic}[1]
\If{$c_0 = \pm\infty$}
  \Return
\EndIf
\State $(c_0, c_1) \gets \textsc{QuickTwoSum}(c_0, c_1)$
\State $s_0 \gets c_0;\; s_1 \gets c_1$
\If{$s_1 \ne 0$}
  \State $(s_1, s_2) \gets \textsc{QuickTwoSum}(s_1, c_2)$
\Else
  \State $(s_0, s_1) \gets \textsc{QuickTwoSum}(s_0, c_2)$
\EndIf
\State $(c_0, c_1, c_2) \gets (s_0, s_1, s_2)$
\end{algorithmic}

\subsection{\texttt{renorm(c0, c1, c2, c3)} --- 4-component}

This variant is needed after TD$+$double and TD$\times$double
operations.  It first performs a bottom-up sweep of
\textsc{QuickTwoSum} to collect carries, then a top-down sweep to
compress into three non-overlapping components:
\begin{algorithmic}[1]
\State \textit{// Bottom-up sweep}
\State $(s, c_3) \gets \textsc{QuickTwoSum}(c_2, c_3)$
\State $(s, c_2) \gets \textsc{QuickTwoSum}(c_1, s)$
\State $(c_0, c_1) \gets \textsc{QuickTwoSum}(c_0, s)$
\State \textit{// Top-down compression with zero-skipping}
\State $s_0 \gets c_0;\; s_1 \gets c_1$
\If{$s_1 \ne 0$}
  \State $(s_1, s_2) \gets \textsc{QuickTwoSum}(s_1, c_2)$
  \State $s_2 \gets s_2 + c_3$ \Comment{absorb remainder}
\Else
  \State $(s_0, s_1) \gets \textsc{QuickTwoSum}(s_0, c_2)$
  \State \textit{continue similarly with $c_3$}
\EndIf
\State $(c_0, c_1, c_2) \gets (s_0, s_1, s_2)$
\end{algorithmic}

\subsection{\texttt{renorm(c0, c1, c2, c3, c4)} --- 5-component}

Required after TD$\times$TD multiplication, which produces five
significant partial sums.  The structure is analogous: a 4-step
bottom-up sweep followed by a top-down compression with
zero-skipping at each level.  The final $c_3$ value (and beyond)
is discarded, truncating to triple-double precision.

Table~\ref{tab:renorm-costs} summarizes the renormalization costs.
All consist entirely of \textsc{QuickTwoSum} calls (3~flops each),
so FMA availability is irrelevant.

\begin{table}[h]
\centering
\caption{Renormalization costs (worst-case branch path).}
\label{tab:renorm-costs}
\begin{tabular}{lcl}
\toprule
Variant & Flops & Composition \\
\midrule
\texttt{td::renorm}$(c_0,c_1,c_2)$             & 6  & $2\times$QTS \\
\texttt{td::renorm}$(c_0,c_1,c_2,c_3)$         & 15 & $5\times$QTS \\
\texttt{td::renorm}$(c_0,c_1,c_2,c_3,c_4)$     & 21 & $7\times$QTS \\
\texttt{qd::renorm}$(c_0,c_1,c_2,c_3)$         & 15 & $5\times$QTS \\
\texttt{qd::renorm}$(c_0,c_1,c_2,c_3,c_4)$     & 21 & $7\times$QTS \\
\bottomrule
\end{tabular}
\end{table}


\section{Addition}
\label{sec:add}

\subsection{TD $+$ double}

Adding a scalar $b$ to a triple-double $a = (a_0, a_1, a_2)$ uses
a simple carry-propagation chain of \textsc{TwoSum} operations:
\begin{algorithmic}[1]
\State $(c_0, e) \gets \textsc{TwoSum}(a_0, b)$
\State $(c_1, e) \gets \textsc{TwoSum}(a_1, e)$
\State $(c_2, e) \gets \textsc{TwoSum}(a_2, e)$
\State \textsc{Renorm4}$(c_0, c_1, c_2, e)$
\State \Return $(c_0, c_1, c_2)$
\end{algorithmic}
Cost: 3 \textsc{TwoSum} + 1 renormalization.

\subsection{TD $+$ TD: IEEE-style merge}

The addition of two triple-doubles uses a merge-based algorithm
analogous to the QD library's \texttt{ieee\_add}.  The six components
from $a$ and $b$ are merged by descending magnitude:
\begin{algorithmic}[1]
\State $i \gets 0;\; j \gets 0;\; k \gets 0$
\State Pick the larger of $|a_i|$ and $|b_j|$ as $u$; advance the
  chosen index
\State Pick the next largest as $v$; advance
\State $(u, v) \gets \textsc{QuickTwoSum}(u, v)$
\While{$k < 3$}
  \State Pick next largest component $t$ from remaining $a_i$ or $b_j$
  \State $s \gets \textsc{QuickThreeAccum}(u, v, t)$
  \If{$s \ne 0$} $x_k \gets s;\; k \gets k + 1$ \EndIf
\EndWhile
\State Accumulate any leftover components into $x_3$
\State \textsc{Renorm4}$(x_0, x_1, x_2, x_3)$
\State \Return $(x_0, x_1, x_2)$
\end{algorithmic}
Here \textsc{QuickThreeAccum} absorbs a new value into the running
pair $(u, v)$ and emits a fully resolved high-order word when
possible.

\subsection{TD $-$ TD: direct differencing}

Subtraction of two triple-doubles uses component-wise
\textsc{TwoDiff}, followed by carry propagation via
\textsc{TwoSum} and \textsc{ThreeSum}:
\begin{algorithmic}[1]
\State $(s_0, t_0) \gets \textsc{TwoDiff}(a_0, b_0)$
\State $(s_1, t_1) \gets \textsc{TwoDiff}(a_1, b_1)$
\State $(s_2, t_2) \gets \textsc{TwoDiff}(a_2, b_2)$
\State $(s_1, t_0) \gets \textsc{TwoSum}(s_1, t_0)$
\State \textsc{ThreeSum}$(s_2, t_0, t_1)$
\State $t_0 \gets t_0 + t_2$
\State \textsc{Renorm4}$(s_0, s_1, s_2, t_0)$
\State \Return $(s_0, s_1, s_2)$
\end{algorithmic}
This is more efficient than negating $b$ and using
\texttt{ieee\_add}, because it avoids the magnitude-comparison
branches.


\section{Multiplication}
\label{sec:mul}

\subsection{TD $\times$ double}

Multiplying $(a_0, a_1, a_2)$ by a scalar $b$ requires three
\textsc{TwoProd} calls and carry propagation:
\begin{algorithmic}[1]
\State $(p_0, q_0) \gets \textsc{TwoProd}(a_0, b)$
\State $(p_1, q_1) \gets \textsc{TwoProd}(a_1, b)$
\State $(p_2, q_2) \gets \textsc{TwoProd}(a_2, b)$
\State $(p_1, q_0) \gets \textsc{TwoSum}(p_1, q_0)$
\State $(p_2, q_1) \gets \textsc{TwoSum}(p_2, q_1)$
\State $q_0 \gets q_0 + q_1 + q_2$
\State \textsc{Renorm4}$(p_0, p_1, p_2, q_0)$
\State \Return $(p_0, p_1, p_2)$
\end{algorithmic}
Cost: 3 \textsc{TwoProd} + 2 \textsc{TwoSum} + 1 renormalization.

When FMA is available, \textsc{TwoProd} reduces to two operations
(one FMA + one subtraction), making the entire multiplication very
efficient.

\subsection{TD $\times$ TD}

The product of $a = (a_0, a_1, a_2)$ and $b = (b_0, b_1, b_2)$
generates nine cross-products $a_i \cdot b_j$.  Of these, only those
contributing to the first $\approx 159$ bits of the result need to be
computed with full error tracking.  The algorithm computes six
\textsc{TwoProd} calls for the terms of order $\le 2$ (i.e.,
$i + j \le 2$), then accumulates higher-order terms approximately:
\begin{algorithmic}[1]
\State $(p_0, q_0) \gets \textsc{TwoProd}(a_0, b_0)$ \Comment{order 0}
\State $(p_1, q_1) \gets \textsc{TwoProd}(a_0, b_1)$ \Comment{order 1}
\State $(p_2, q_2) \gets \textsc{TwoProd}(a_1, b_0)$ \Comment{order 1}
\State $(p_3, q_3) \gets \textsc{TwoProd}(a_0, b_2)$ \Comment{order 2}
\State $(p_4, q_4) \gets \textsc{TwoProd}(a_1, b_1)$ \Comment{order 2}
\State $(p_5, q_5) \gets \textsc{TwoProd}(a_2, b_0)$ \Comment{order 2}
\State \textit{// Accumulate order-1 terms}
\State \textsc{ThreeSum}$(p_1, p_2, q_0)$
\State \textsc{ThreeSum}$(p_2, q_1, q_2)$
\State \textsc{ThreeSum}$(p_3, p_4, p_5)$
\State \textit{// Merge order-2 contributions}
\State $(s_0, t_0) \gets \textsc{TwoSum}(p_2, p_3)$
\State $(s_1, t_1) \gets \textsc{TwoSum}(q_1, p_4)$
\State $s_2 \gets q_2 + p_5$
\State $(s_1, t_0) \gets \textsc{TwoSum}(s_1, t_0)$
\State $s_2 \gets s_2 + t_0 + t_1$
\State \textit{// Approximate higher-order terms}
\State $s_1 \gets s_1 + a_1 b_2 + a_2 b_1 + q_0 + q_3 + q_4 + q_5$
\State \textsc{Renorm5}$(p_0, p_1, s_0, s_1, s_2)$
\State \Return $(p_0, p_1, s_0)$
\end{algorithmic}
The key insight is that the three order-2 cross-products and their
error terms contribute to the third component, while the even
higher-order terms ($a_1 b_2$, $a_2 b_1$, and the six error words
$q_0, \ldots, q_5$) are accumulated in standard double arithmetic
because they affect only the lowest bits.

\subsection{Squaring}

Currently, squaring is implemented as \texttt{sqr(a) = a * a}.
A dedicated squaring routine exploiting the symmetry
$a_i b_j = a_j b_i$ could reduce the \textsc{TwoProd} count
from 6 to 4, but this optimization is deferred.

\begin{remark}[Multiplication has no overflow protection]
\label{rem:mul-no-rescale}
Unlike division and square root (which apply explicit rescaling
when $|a_0| > 2^{969}$), the multiplication operators
TD$\times$double and TD$\times$TD do \emph{not} include overflow
guards.  They rely entirely on the \textsc{Split} function's
built-in prescaling at $|a| > 2^{996}$
(Remark~\ref{rem:split}).  This is safe because multiplication
does not amplify the exponent of a single operand's components:
the products $a_i \cdot b_j$ have exponents bounded by
$\exp(a_i) + \exp(b_j)$, and if both inputs are representable
doubles, all partial products remain below $2^{1023}$.  The
\textsc{Split} threshold at $2^{996}$ provides 27 bits of
headroom for the $\sigma = 2^{27} + 1$ multiplier.

Division and square root, by contrast, generate intermediate
residuals and quotient digits whose exponent relationships are
less predictable, necessitating the additional $2^{969}$ threshold
(see Remark~\ref{rem:threshold969} below).
\end{remark}


\section{Division}
\label{sec:div}

\subsection{TD $\div$ double}

Division by a scalar uses the standard long-division approach with
three quotient digits:
\begin{algorithmic}[1]
\Require $b \ne 0$
\State \textit{// Overflow-safe rescaling}
\If{$|a_0| > 2^{969}$}
  \State $a \gets a \cdot 2^{-53}$;\; \texttt{rescale} $\gets$ true
\EndIf
\State $q_0 \gets a_0 / b$
\State $(t_0, t_1) \gets \textsc{TwoProd}(q_0, b)$
\State $r \gets a - (t_0, t_1, 0)$
\State $q_1 \gets r_0 / b$
\State $(t_0, t_1) \gets \textsc{TwoProd}(q_1, b)$
\State $r \gets r - (t_0, t_1, 0)$
\State $q_2 \gets r_0 / b$
\State \textsc{Renorm3}$(q_0, q_1, q_2)$
\If{\texttt{rescale}}
  \State result $\gets$ result $\cdot 2^{53}$
\EndIf
\State \Return $(q_0, q_1, q_2)$
\end{algorithmic}

\begin{remark}[Overflow protection in division]
\label{rem:threshold969}
The threshold $2^{969}$ and the rescaling factor $2^{-53}$ are
derived from the \textsc{Split} threshold $2^{996}$
(Remark~\ref{rem:split}).

In the TD$\div$TD algorithm, the residual step computes
$r = a' - b \cdot q_0$ where $q_0 = a'_0 / b_0$.  The product
$b \cdot q_0$ is a TD$\times$double operation, which internally
calls \textsc{TwoProd}$(b_i, q_0)$ for each component $b_i$.
Inside the non-FMA \textsc{TwoProd}, the \textsc{Split} function
computes $\sigma \cdot b_i$ where $\sigma = 2^{27} + 1$.  If the
divisor $b$ has not been rescaled, $|b_0|$ can be as large as
$\approx 2^{1023}$, and \textsc{Split} handles this via its
built-in prescaling at $2^{996}$.

The danger arises from the \emph{quotient digit} $q_0$.  If
$|a_0|$ is large and $|b_0|$ is moderate, $q_0 \approx a_0 / b_0$
inherits the large exponent of $a_0$.  \textsc{TwoProd}$(q_0, b)$
then requires \textsc{Split}$(q_0)$, which computes
$\sigma \cdot q_0$.  For this to stay below $2^{1023}$, we need
$|q_0| < 2^{996}$.  After rescaling $a$ by $2^{-53}$, we have
$|q_0| \le |a_0| \cdot 2^{-53} / |b_0| \le 2^{1023-53} / 1 =
2^{970}$, which is safely below $2^{996}$.

The choice $2^{969}$ as the trigger (rather than $2^{996}$ directly)
provides a conservative margin: $996 - 969 = 27$ bits, matching
the $\log_2 \sigma = 27$ bits consumed by the splitting constant.

Note that the rescaling is applied \emph{once}, not iteratively.
After a single rescaling by $2^{-53}$, the worst-case numerator
exponent is $2^{970}$, which exceeds the trigger $2^{969}$ but
remains below the actual danger zone $2^{996}$.  The trigger
is a conservative pre-check, not a precise boundary; the true
safety invariant is $|q_i| < 2^{996}$, which the single
rescaling guarantees.
\end{remark}

\subsection{TD $\div$ TD}

The algorithm extracts quotient digits by Newton-style residual
correction, with each refinement step using a full TD$\times$double
multiplication and a full TD$-$TD subtraction.  Two variants are
provided, distinguished only by the number of correction iterations:
\begin{algorithmic}[1]
\Require $b \ne 0$
\State Apply rescaling if $|a_0| > 2^{969}$
\State $q_0 \gets a_0 / b_0$
\State $r \gets a - b \cdot q_0$ \Comment{TD $\times$ double, TD $-$ TD}
\State $q_1 \gets r_0 / b_0$
\State $r \gets r - b \cdot q_1$
\State $q_2 \gets r_0 / b_0$
\If{accurate variant}
  \State $r \gets r - b \cdot q_2$
  \State $q_3 \gets r_0 / b_0$
  \State \textsc{Renorm4to3}$(q_0, q_1, q_2, q_3)$
\Else
  \State \textsc{Renorm3}$(q_0, q_1, q_2)$
\EndIf
\State Apply inverse rescaling if needed
\State \Return $(q_0, q_1, q_2)$
\end{algorithmic}
Three quotient digits already saturate the three-component
representation: each iteration captures $\approx 53$ new bits, and
three iterations yield $\approx 159$ bits, matching the
$\approx \varepsilon^3$ relative accuracy target.  The accurate
variant performs a fourth scalar division and one additional
TD$-$TD subtraction, giving roughly one extra guard bit in the
least-significant component---valuable when the result is then fed
into a long chain of further operations.

This sloppy/accurate split mirrors the
\texttt{qd\_real::\{sloppy,accurate\}\_div} convention, where sloppy
runs four iterations and accurate runs five.  Notably, an earlier
TD implementation tried a different notion of ``sloppy'' that
replaced the residual TD$-$TD subtraction with a cheap
\textsc{TwoDiff}-plus-scalar chain.  This is unsound at three
components: collapsing the residual to two scalars discards exactly
the information needed to extract $q_2$ at $O(\varepsilon^3)$, and
the result silently degrades to DD precision.  The current
implementation therefore uses full TD$-$TD subtractions in both
variants.


\section{Computational Cost Analysis}
\label{sec:cost}

Table~\ref{tab:flop-main} gives the total binary64 flop count for
each same-type addition, subtraction, multiplication, and division,
derived by composing the primitive costs from
Tables~\ref{tab:eft-costs} and~\ref{tab:renorm-costs}.
Addition and subtraction use only
\textsc{TwoSum}/\textsc{TwoDiff}/\textsc{QuickTwoSum}
(all independent of FMA), so FMA affects only the multiplication
and division columns.
The default variant for each type is marked with~$\star$.

\begin{table}[ht]
\centering
\caption{Binary64 flop counts for DD, TD, and QD arithmetic.
  \#TP = number of \textsc{TwoProd} calls.
  FMA counts each fused multiply-add as 1~flop.
  Division costs exclude 2--3 scalar \texttt{double/double}
  divisions (counted separately as ``div'' in the Breakdown column).
  $\star$ = default variant.}
\label{tab:flop-main}
\begin{tabular}{llrrrl}
\toprule
Operation & Variant & FMA & No FMA & \#TP & Breakdown \\
\midrule
DD$+$DD & sloppy$\,\star$ & 11 & 11 & 0
  & $1\,\text{TS}+2+1\,\text{QTS}$ \\
DD$+$DD & ieee & 20 & 20 & 0
  & $2\,\text{TS}+2+2\,\text{QTS}$ \\
DD$\times$DD & ---$\,\star$ & 9 & 24 & 1
  & $1\,\text{TP}+2\,\text{mul}+2\,\text{add}+1\,\text{QTS}$ \\
\midrule
TD$+$TD & sloppy$\,\star$ & 41 & 41 & 0
  & $4\,\text{TS}+2+\text{renorm4}$ \\
TD$+$TD & ieee & 55--68 & 55--68 & 0
  & $1\,\text{QTS}+3\text{--}4\,\text{QTA}+\text{renorm4}$ \\
TD$-$TD & sloppy$\,\star$ & 41 & 41 & 0
  & $3\,\text{TD}+1\,\text{TS}+2+\text{renorm4}$ \\
TD$-$TD & accurate & 58 & 58 & 0
  & $3\,\text{TD}+1\,\text{TS}+1\,\text{3S}+1+\text{renorm4}$ \\
TD$\times$TD & ---$\,\star$ & 116 & 206 & 6
  & $6\,\text{TP}+3\,\text{3S}+3\,\text{TS}+8+\text{renorm5}$ \\
TD$\div$TD & accurate$\,\star$ & $4\text{d}+240$ & $4\text{d}+375$ & 9
  & $3\times(\text{TD}{\times}\text{d}+\text{TD}{\text{-}}\text{TD})+\text{rn4}$ \\
TD$\div$TD & sloppy & $3\text{d}+158$ & $3\text{d}+248$ & 6
  & $2\times(\text{TD}{\times}\text{d}+\text{TD}{\text{-}}\text{TD})+\text{rn3}$ \\
\midrule
QD$+$QD & sloppy$\,\star$ & 84 & 84 & 0
  & $5\,\text{TS}+1\,\text{3S}+1\,\text{3S2}+2+\text{renorm5}$ \\
QD$\times$QD & sloppy & 120 & 210 & 6
  & $6\,\text{TP}+3\,\text{3S}+3\,\text{TS}+12+\text{renorm5}$ \\
QD$\times$QD & accurate$\,\star$ & 182 & 332 & 10
  & $10\,\text{TP}+3\,\text{3S}+11\,\text{TS}+11+\text{renorm5}$ \\
\bottomrule
\end{tabular}
\\[2pt]
\footnotesize
TS = \textsc{TwoSum}(6), TD = \textsc{TwoDiff}(6),
QTS = \textsc{QuickTwoSum}(3), TP = \textsc{TwoProd}(2/17),
3S = \textsc{ThreeSum}(18), 3S2 = \textsc{ThreeSum2}(13),
QTA = \textsc{QuickThreeAccum}(12),
d = scalar \texttt{double} division,
rn3/rn4 = \texttt{renorm} 3/4-component.
Default variants: addition/subtraction = sloppy
(\texttt{QD\_IEEE\_ADD} selects ieee/accurate);
division = accurate
(\texttt{QD\_SLOPPY\_DIV} selects sloppy).
Division costs assume default (sloppy) subtraction inside
\texttt{accurate\_div}.
\end{table}

Several observations are noteworthy:

\begin{itemize}
\item The TD \texttt{sloppy\_add} (41~flops) replaces the
  \textsc{ThreeSum} in the carry-propagation stage with plain
  double-precision addition, saving 16~flops vs.\ the accurate
  variant (58~flops for subtraction, comparable for addition) while
  maintaining Cray-style error bounds.  This matches the DD/QD
  convention where the default addition is the sloppy variant.

\item The TD \texttt{ieee\_add} (55--68~flops) remains available via
  \texttt{QD\_IEEE\_ADD} for applications requiring a tighter IEEE-style
  error bound.  Its cost is data-dependent because the merge loop
  emits a resolved word only when the running pair $(u, v)$ in
  \textsc{QuickThreeAccum} is fully occupied.

\item TD \texttt{sloppy\_div} and \texttt{accurate\_div} differ in the
  \emph{number of correction iterations}, not in the cost of each
  residual: both perform full TD$-$TD subtractions when forming
  $r \gets r - b\cdot q_k$, mirroring the
  \texttt{qd\_real::\{sloppy,accurate\}\_div} convention.
  \texttt{sloppy\_div} runs three iterations ($q_0,q_1,q_2$) and
  delivers $\approx \varepsilon^3$ relative accuracy, which is the
  minimum needed to fill the three TD components;
  \texttt{accurate\_div} runs a fourth iteration ($q_3$) before the
  final \textsc{Renorm4}$\to$3, gaining $\approx 1$~bit of guard
  precision in the least-significant component.  An earlier attempt
  to make \texttt{sloppy\_div} cheaper by replacing the residual
  TD$-$TD subtraction with a \textsc{TwoDiff} + scalar chain was
  abandoned: collapsing the residual to two scalars discards the
  information needed to extract $q_2$ at $O(\varepsilon^3)$, and the
  result degrades to DD precision.  The cost of the extra iteration
  is one additional TD$\times$double, one TD$-$TD subtraction, one
  scalar division, and \textsc{Renorm4}$\to$3 in place of
  \textsc{Renorm3}.

\item TD$\times$TD and QD$\times$QD \texttt{sloppy\_mul} have the
  same structure (6 \textsc{TwoProd} + 3 \textsc{ThreeSum}), but
  QD accumulates 4 additional approximate multiply-adds for the
  $O(\varepsilon^3)$ terms.

\item FMA provides a $\approx 1.8\times$ speedup for multiplication
  (e.g., 116 vs.\ 206 for TD$\times$TD) and a $\approx 1.6\times$
  speedup for division, but does not affect addition or subtraction.
\end{itemize}

Table~\ref{tab:flop-ratio} gives the normalized cost ratios using
default variants.  TD achieves approximately 74\% of QD's
precision (157 vs.\ 212 bits) at approximately 49\% of QD's
addition cost and 64\% of QD's multiplication cost, confirming
that it fills the DD--QD gap cost-effectively.

\begin{table}[h]
\centering
\caption{Flop ratios using default variants.  ``Ratio'' is
  relative to DD.  Division excludes scalar \texttt{double/double}
  divisions.}
\label{tab:flop-ratio}
\begin{tabular}{lccccc}
\toprule
      & & Addition & Multiplication & \multicolumn{2}{c}{TD / QD} \\
\cmidrule(lr){5-6}
Type  & Bits & Flops (Ratio) & FMA / No FMA (Ratio) & Add & Mul \\
\midrule
DD    & 106  & 11 ($1.0\times$) & 9 / 24 ($1.0\times$)
  & --- & --- \\
TD    & 157  & 41 ($3.7\times$) & 116 / 206 ($12.9\times$ / $8.6\times$)
  & 49\% & 64\% / 62\% \\
QD    & 212  & 84 ($7.6\times$) & 182 / 332 ($20.2\times$ / $13.8\times$)
  & --- & --- \\
\bottomrule
\end{tabular}
\\[2pt]
\footnotesize
Default variants: DD/TD/QD addition = sloppy,
TD multiplication = (only one variant),
QD multiplication = accurate.
TD/QD column: addition $41/84 = 49\%$;
multiplication $116/182 = 64\%$ (FMA), $206/332 = 62\%$ (no FMA).
\end{table}


\section{Type Conversion and Mixed-Precision Arithmetic}
\label{sec:mixed}

A key design goal of \texttt{td\_real} is seamless interoperability
with the existing \texttt{dd\_real} and \texttt{qd\_real} types.
This section describes the conversion algorithms and the
mixed-precision operator semantics.

\subsection{Conversion between TD and QD}

\subsubsection{QD $\to$ TD: truncation with renormalization}

A quad-double $q = (q_0, q_1, q_2, q_3)$ has four components; to
convert it to a triple-double we must discard the lowest component
while preserving the non-overlapping invariant.  Simply dropping
$q_3$ would leave a valid triple-double only if $(q_0, q_1, q_2)$
already satisfies the invariant, which is not guaranteed after
QD arithmetic.  The implementation therefore applies a 4-component
renormalization before truncation:
\begin{algorithmic}[1]
\Function{FromQdTruncate}{$q$}
  \State $(x_0, x_1, x_2, x_3) \gets (q_0, q_1, q_2, q_3)$
  \State \textsc{Renorm4}$(x_0, x_1, x_2, x_3)$
    \Comment{restores non-overlapping}
  \State \Return $(x_0, x_1, x_2)$
    \Comment{$x_3$ discarded}
\EndFunction
\end{algorithmic}
The renormalization step ensures that the carry from $x_3$ is
properly propagated into $x_2$, $x_1$, and $x_0$ before the
truncation.  The relative error of this conversion is bounded by
$\varepsilon_{\mathrm{td}} = 2^{-157}$.

This function is used by the \texttt{explicit} constructor
\texttt{td\_real(const qd\_real \&qd)} and the assignment operator
\texttt{td\_real::operator=(const qd\_real \&)}, as well as
internally by all mixed TD--QD compound-assignment operators.

\subsubsection{TD $\to$ QD: zero-extension}

The reverse conversion is lossless: a triple-double
$(a_0, a_1, a_2)$ is promoted to a quad-double by appending a
zero fourth component:
\begin{algorithmic}[1]
\Function{ToQdConversion}{$a$}
  \State \Return $(a_0, a_1, a_2, 0)$
\EndFunction
\end{algorithmic}
Since the non-overlapping invariant of the original triple-double
is a subset of the quad-double invariant, no renormalization is
needed.  The result is exact.

\subsection{Conversion between TD and DD}

\subsubsection{DD $\to$ TD}

A double-double $(d_0, d_1)$ is promoted to triple-double by
zero-extending the third component:
\begin{lstlisting}
td_real(const dd_real &dd) {
  x[0] = dd._hi();
  x[1] = dd._lo();
  x[2] = 0.0;
}
\end{lstlisting}
This is exact and preserves the non-overlapping invariant, since
the DD invariant $|d_1| \le \frac{1}{2}\,\mathrm{ulp}(d_0)$ is
strictly stronger than what TD requires for its first two components.

\subsubsection{TD $\to$ DD: lossy truncation}

Truncating from triple-double to double-double discards $a_2$ but
must account for its contribution to the DD pair:
\begin{lstlisting}
dd_real to_dd_real(const td_real &a) {
  dd_real result(a[0], a[1]);
  result += a[2];  // absorb low component
  return result;
}
\end{lstlisting}
The \texttt{+=} operation on \texttt{dd\_real} re-normalizes the
pair, ensuring that the carry from $a_2$ is correctly propagated.
The relative error is $O(\varepsilon_{\mathrm{dd}}) = O(2^{-106})$.

\subsection{Mixed-precision arithmetic operators}

The C++ operator overloads follow a \emph{promotion rule}:

\begin{center}
\begin{tabular}{lccl}
\toprule
Operation & Left type & Right type & Result type \\
\midrule
\texttt{td + d}   & TD     & double & TD \\
\texttt{td + dd}  & TD     & DD     & TD \\
\texttt{td + td}  & TD     & TD     & TD \\
\texttt{td + qd}  & TD     & QD     & QD \\
\texttt{qd + td}  & QD     & TD     & QD \\
\bottomrule
\end{tabular}
\end{center}

The general principle is:
\begin{itemize}
\item When \emph{both} operands fit within TD precision (double,
  DD, or TD), the result is TD.  The narrower operand is promoted
  to TD before the operation.
\item When \emph{either} operand is QD, the result is QD.  The TD
  operand is promoted to QD via \textsc{ToQdConversion}, and the
  operation is performed using the existing QD arithmetic.
\end{itemize}

The same promotion rule applies to subtraction, multiplication,
and division.  This design keeps TD as a ``closed'' type for
operations among double/DD/TD, while automatically widening to QD
when QD precision is requested by either operand.

\subsubsection{Implementation of TD $\otimes$ QD}

For all four binary operators $\otimes \in \{+, -, \times, /\}$,
the TD--QD mixed operation is:
\begin{lstlisting}
qd_real operator+(const td_real &a, const qd_real &b) {
  return td::to_qd_conversion(a) + b;  // QD + QD
}
\end{lstlisting}
The conversion \textsc{ToQdConversion} is $O(1)$ (a simple copy
with a zero append), so the cost is dominated by the QD operation
itself.

\subsubsection{Implementation of TD $\otimes$ DD}

For DD operands, promotion to TD is used:
\begin{lstlisting}
td_real operator+(const td_real &a, const dd_real &b) {
  return a + td_real(b);  // TD + TD
}
\end{lstlisting}
The DD-to-TD promotion is also $O(1)$ (copy two doubles, set third
to zero).

\subsection{Compound-assignment operators with QD}

The compound-assignment operators (\texttt{+=}, \texttt{-=},
\texttt{*=}, \texttt{/=}) for QD right-hand sides require special
care because the result type of TD $\otimes$ QD is QD, but the
target variable is TD.  The implementation therefore performs the
operation in QD and then truncates:
\begin{lstlisting}
td_real &td_real::operator+=(const qd_real &a) {
  *this = td::from_qd_truncate(
            td::to_qd_conversion(*this) + a);
  return *this;
}
\end{lstlisting}
The sequence is: (1)~promote \texttt{*this} to QD, (2)~perform QD
addition, (3)~truncate the QD result back to TD via
\textsc{FromQdTruncate}.  This ensures that the full QD precision
of the intermediate result is correctly rounded into the TD target.

The same pattern applies to \texttt{-=}, \texttt{*=}, and
\texttt{/=} with QD operands.

\subsection{Cross-type comparison}

Comparison operators between TD and QD are implemented by promoting
the TD operand to QD:
\begin{lstlisting}
bool operator<(const td_real &a, const qd_real &b) {
  return td::to_qd_conversion(a) < b;
}
\end{lstlisting}
This is exact because TD$\to$QD is lossless.  For TD vs.\ DD
comparisons, the DD operand is promoted to TD, which is also exact
(DD$\to$TD is lossless).

The C API provides explicit comparison functions for mixed types:
\texttt{c\_td\_comp\_td\_qd} and \texttt{c\_td\_comp\_qd\_td},
which convert to TD before comparing (truncating the QD operand).
This is a deliberate design choice: the C API operates entirely
in TD space, so QD arguments are truncated to TD precision before
comparison.  This differs from the C++ operator semantics where
the comparison is performed in QD space.

\subsection{Interaction with DD's compound-assignment operators}

The \texttt{dd\_real} type is also extended with compound-assignment
operators that accept \texttt{td\_real}:
\begin{lstlisting}
dd_real &dd_real::operator+=(const td_real &a) {
  *this += to_dd_real(a);
  return *this;
}
\end{lstlisting}
Here the TD operand is truncated to DD before the DD addition.
Similarly, \texttt{qd\_real} gains compound-assignment operators
for TD operands by promoting TD to QD:
\begin{lstlisting}
qd_real &qd_real::operator+=(const td_real &a) {
  *this += td::to_qd_conversion(a);
  return *this;
}
\end{lstlisting}
These extensions make all three types (DD, TD, QD) mutually
interoperable in both directions.

\subsection{Structural modifications to the DD and QD headers}

Integrating \texttt{td\_real} into the QD library required
coordinated changes to the existing \texttt{dd\_real.h} and
\texttt{qd\_real.h} headers.  These modifications are worth
documenting because they define the ``surface area'' of the TD
extension visible to code that uses DD or QD without directly
including \texttt{td\_real.h}.

\subsubsection{Changes to \texttt{dd\_real.h}}

A forward declaration \texttt{struct td\_real;} is added before
the \texttt{dd\_real} struct definition.  The following TD-aware
members and free functions are added:
\begin{itemize}
\item \textbf{Compound-assignment operators}:
  \texttt{operator+=}, \texttt{-=}, \texttt{*=}, \texttt{/=}
  accepting \texttt{const td\_real \&}.
\item \textbf{Binary arithmetic operators}:
  \texttt{operator+}, \texttt{-}, \texttt{*}, \texttt{/} for
  (DD, TD) and (TD, DD) pairs, all returning \texttt{td\_real}.
\item \textbf{Comparison operators}:
  \texttt{==}, \texttt{!=}, \texttt{<}, \texttt{>}, \texttt{<=},
  \texttt{>=} for (DD, TD) and (TD, DD) pairs.
\end{itemize}
These additions follow the promotion rule: DD $\otimes$ TD yields
TD, never DD.

\subsubsection{Changes to \texttt{qd\_real.h}}

Similarly, \texttt{qd\_real.h} gains a forward declaration and:
\begin{itemize}
\item \textbf{Compound-assignment operators}:
  \texttt{operator+=}, \texttt{-=}, \texttt{*=}, \texttt{/=}
  accepting \texttt{const td\_real \&}.
\item \textbf{Binary arithmetic operators}:
  \texttt{operator+}, \texttt{-}, \texttt{*}, \texttt{/} for
  (QD, TD) and (TD, QD) pairs, all returning \texttt{qd\_real}.
\item \textbf{Comparison operators}:
  \texttt{==}, \texttt{!=}, \texttt{<}, \texttt{>}, \texttt{<=},
  \texttt{>=} for (QD, TD) and (TD, QD) pairs.
\item \textbf{C API additions} (\texttt{c\_qd.h}):
  Functions like \texttt{c\_qd\_add\_td\_qd},
  \texttt{c\_qd\_selfadd\_td}, \texttt{c\_qd\_copy\_td},
  \texttt{c\_qd\_comp\_qd\_td}, etc.
\end{itemize}
These follow the QD promotion rule: QD $\otimes$ TD yields QD.

\subsubsection{Changes to \texttt{c\_dd.h}}

The C wrapper header for DD also gains TD-aware functions:
\texttt{c\_dd\_add\_dd\_td}, \texttt{c\_dd\_sub\_td\_dd},
\texttt{c\_dd\_mul\_dd\_td}, \texttt{c\_dd\_div\_td\_dd},
\texttt{c\_dd\_copy\_td}, \texttt{c\_dd\_selfadd\_td},
\texttt{c\_dd\_comp\_dd\_td}, etc.  These enable C (and Fortran)
code to perform mixed DD--TD arithmetic through the C binding layer.

\subsubsection{Header inclusion order}

The inclusion order matters:
\texttt{dd\_real.h} $\to$ \texttt{qd\_real.h} $\to$
\texttt{td\_real.h}.  Since \texttt{dd\_real.h} and
\texttt{qd\_real.h} only use a forward declaration of
\texttt{td\_real}, they do not depend on the full TD definition.
The inline implementations of the TD-aware operators on DD and QD
are placed in \texttt{td\_inline.h}, which is included at the end
of \texttt{td\_real.h} (guarded by \texttt{QD\_INLINE}).  This
avoids circular dependencies while keeping the operators inlinable.

\subsection{Summary of conversion costs}

\begin{center}
\begin{tabular}{lccc}
\toprule
Conversion & Cost & Exact? & Renorm needed? \\
\midrule
DD $\to$ TD & $O(1)$: copy + zero & Yes & No \\
TD $\to$ DD & $O(1)$: DD \texttt{+=} & No ($\varepsilon_{\mathrm{dd}}$) & DD renorm \\
TD $\to$ QD & $O(1)$: copy + zero & Yes & No \\
QD $\to$ TD & $O(1)$: 4-renorm + truncate & No ($\varepsilon_{\mathrm{td}}$) & Yes (4-comp) \\
\bottomrule
\end{tabular}
\end{center}


\section{Square Root}
\label{sec:sqrt}

The square root uses Newton's method with a double-precision seed:
\begin{algorithmic}[1]
\Require $a \ge 0$
\If{$|a_0| > 2^{969}$} \Comment{overflow-safe rescaling}
  \State \Return $2 \cdot \textsc{Sqrt}(a/4)$
    \Comment{recursive call}
\EndIf
\State $x \gets (\sqrt{a_0}, 0, 0)$ \Comment{double-precision seed}
\State $x \gets x + (a - x^2) / (2x)$ \Comment{first Newton step}
\State $x \gets x + (a - x^2) / (2x)$ \Comment{second Newton step}
\State \Return $x$
\end{algorithmic}
Each Newton step has quadratic convergence, roughly doubling the
number of correct bits per step.  The seed provides $\approx 53$
bits; after two iterations we have $\approx 212$ bits, which
exceeds the 157 bits needed.  The excess is absorbed during
renormalization.

\begin{remark}[Recursive rescaling in square root]
\label{rem:sqrt-recurse}
Unlike division, which applies a single non-recursive rescaling
by $2^{-53}$ (Remark~\ref{rem:threshold969}), the square root
uses \emph{recursive} rescaling: \texttt{sqrt} calls itself with
the argument divided by~4.  The identity used is
$\sqrt{a} = 2\sqrt{a/4}$.

The recursion depth depends on the input exponent.
\textsc{MulPwr2}$(a, 0.25)$ reduces $\exp(a_0)$ by~2 per call.
Starting from the maximum double exponent 1023, the recursion
terminates when $\exp(a_0) \le 969$, requiring at most
$\lceil(1023 - 969)/2\rceil = 27$ recursive calls.  While this
is computationally wasteful for extreme inputs, such values are
rare in practice.

The reason recursive rescaling is needed (rather than a single
large rescale) is that the Newton iteration computes $x^2$ via
the TD$\times$TD multiplication, which in turn calls
\textsc{TwoProd} on the components of~$x$.  If $a_0 \approx
2^{e}$, then $x_0 \approx 2^{e/2}$, and the cross-products
$x_0 \cdot x_1$ in the squaring step have exponent $\approx
e/2 + (e/2 - 53) = e - 53$.  For the \textsc{Split} operation
on these cross-products to remain safe, we need $e - 53 < 996$,
i.e., $e < 1049$, which is always satisfied for representable
doubles.  However, the more stringent constraint comes from
the \emph{division} $/(2x)$ in the Newton step, which triggers
the division rescaling path (Remark~\ref{rem:threshold969}) if
the residual $(a - x^2)$ has $|a_0 - x_0^2| > 2^{969}$.
By reducing the input to $\le 2^{969}$, the square root
ensures that all intermediate values in both the squaring and
the division sub-operations remain in the safe regime.
\end{remark}

For the $n$th root (\texttt{nroot}), the implementation solves
$1/x^n - a = 0$ via Newton iteration in the reciprocal, then
returns $1/x$.  Three Newton steps are used to converge from a
double-precision seed.


\section{Transcendental Functions}
\label{sec:transcendental}

The transcendental functions follow the same general strategies as in
the DD and QD libraries, adapted for triple-double precision.

\subsection{Exponential}

The exponential $e^a$ uses argument reduction and repeated squaring:
\begin{enumerate}
\item \textbf{Range check.}  Handle NaN, $\pm\infty$, zero, and
  overflow/underflow ($|a_0| \ge 709$).
\item \textbf{Argument reduction.}  Compute
  $m = \lfloor a / \ln 2 + 0.5 \rfloor$ and
  $r = (a - m \ln 2) / k$ where $k = 2^{12} = 4096$.
  This makes $|r| \le \ln 2 / (2k) \approx 8.5 \times 10^{-5}$.
\item \textbf{Taylor series.}  Evaluate $s = e^r - 1$ via the
  Taylor series $s = r + r^2/2! + r^3/3! + \cdots$, terminating
  when $|\text{term}| < \varepsilon_{\mathrm{td}} / k$.
\item \textbf{Repeated squaring.}  Recover $e^{a/k^2}$ from $s$ by
  applying $s \gets 2s + s^2$ a total of 12 times
  ($k = 2^{12}$).
\item \textbf{Scaling.}  Return $(s + 1) \cdot 2^m$ via
  \texttt{ldexp}.
\end{enumerate}

\subsection{Logarithm}

The natural logarithm uses Halley-like Newton iteration on $f(x) =
a \cdot e^{-x} - 1$:
\begin{enumerate}
\item \textbf{Initial estimate.}  Use \texttt{std::frexp} to extract
  the exponent $e$ and mantissa $m$, then set
  $x_0 = \ln m + e \cdot \ln 2$ in double precision.
\item \textbf{Newton refinement.}  Apply
  $x \gets x + a \cdot e^{-x} - 1$ three times.  Each step
  roughly triples the number of correct digits (because the
  iteration has cubic convergence for this particular form).
\end{enumerate}

\subsection{Trigonometric functions}

The \texttt{sincos} function (which \texttt{sin}, \texttt{cos},
and \texttt{tan} delegate to) uses a three-stage argument reduction:
\begin{enumerate}
\item \textbf{Full-period reduction.}  Compute
  $z = \mathrm{nint}(a / 2\pi)$ and $r = a - 2\pi z$.
\item \textbf{Quadrant reduction.}  Compute
  $q = \lfloor r / (\pi/2) + 0.5 \rfloor$ and
  $t = r - q \cdot \pi/2$.  The integer $j = q \bmod 4$
  identifies the quadrant.
\item \textbf{Fine reduction.}  Compute
  $q' = \lfloor t / (\pi/16) + 0.5 \rfloor$ and
  $t \gets t - q' \cdot \pi/16$.  The integer $k = q'$
  (with $|k| \le 4$) indexes a table of precomputed
  $(\sin(k\pi/16), \cos(k\pi/16))$ values.
\end{enumerate}

After reduction, $|t| \le \pi/32 \approx 0.098$, and the Taylor
series for $\sin t$ and $\cos t$ converge rapidly.  The
implementation computes $\sin t$ via Taylor and obtains $\cos t$
as $\sqrt{1 - \sin^2 t}$.  The final result is reconstructed using
the angle-addition formula:
\[
  \sin(k\pi/16 + t) = \sin(k\pi/16)\cos t + \cos(k\pi/16)\sin t,
\]
\[
  \cos(k\pi/16 + t) = \cos(k\pi/16)\cos t - \sin(k\pi/16)\sin t,
\]
followed by quadrant adjustment based on $j$.

The table values for $k = 1, 2, 3, 4$ are stored as
triple-double constants (arrays of three doubles) to avoid
precision loss.

\subsection{Inverse trigonometric functions}

\begin{itemize}
\item \texttt{asin}$(a)$ $=$ \texttt{atan2}$(a, \sqrt{1 - a^2})$.
\item \texttt{acos}$(a)$ $=$ \texttt{atan2}$(\sqrt{1 - a^2}, a)$.
\item \texttt{atan}$(a)$ $=$ \texttt{atan2}$(a, 1)$.
\end{itemize}

The core \texttt{atan2}$(y, x)$ function uses Newton iteration.
After normalization $(x, y) \gets (x, y) / \sqrt{x^2 + y^2}$,
a double-precision seed $z_0 = \text{atan2}(y_0, x_0)$ is refined
by three Newton steps of the form:
\[
  z \gets z + \frac{y' - \sin z}{\cos z}
  \quad \text{or} \quad
  z \gets z - \frac{x' - \cos z}{\sin z},
\]
choosing whichever denominator is larger in magnitude.

\subsection{Hyperbolic functions}

\begin{itemize}
\item For $|a| > 0.05$: $\sinh a = (e^a - e^{-a})/2$,
  $\cosh a = (e^a + e^{-a})/2$.
\item For $|a| \le 0.05$: a Taylor series is used for $\sinh$ to
  avoid catastrophic cancellation:
  $\sinh a = a + a^3/3! + a^5/5! + \cdots$.
  Then $\cosh a = \sqrt{1 + \sinh^2 a}$.
\item $\tanh a = \sinh a / \cosh a$, with the same small-argument
  branch.
\end{itemize}

\subsection{Inverse hyperbolic functions}

These use standard identities:
\begin{align}
  \mathrm{asinh}(a) &= \ln(a + \sqrt{a^2 + 1}), \\
  \mathrm{acosh}(a) &= \ln(a + \sqrt{a^2 - 1}), \quad a \ge 1, \\
  \mathrm{atanh}(a) &= \tfrac{1}{2}\ln\!\bigl(\tfrac{1+a}{1-a}\bigr),
    \quad |a| < 1.
\end{align}

\subsection{Integer power and general power}

\texttt{npwr}$(a, n)$ uses binary exponentiation (repeated squaring)
for $O(\log n)$ multiplications.  \texttt{pow}$(a, b) = e^{b \ln a}$.


\section{Precomputed Constants}
\label{sec:constants}

The following constants are stored as \texttt{static const td\_real}
members, each given as three double-precision hex or decimal literals
computed to full triple-double accuracy:

\begin{center}
\begin{tabular}{ll}
\toprule
Constant & Value \\
\midrule
$\pi$          & \texttt{(3.14159265\ldots e+00, 1.22464679\ldots e$-$16, $-$2.99476980\ldots e$-$33)} \\
$2\pi$         & \texttt{(6.28318530\ldots e+00, 2.44929359\ldots e$-$16, $-$5.98953961\ldots e$-$33)} \\
$e$            & \texttt{(2.71828182\ldots e+00, 1.44564689\ldots e$-$16, $-$2.12771710\ldots e$-$33)} \\
$\ln 2$        & \texttt{(6.93147180\ldots e$-$01, 2.31904681\ldots e$-$17, 5.70770843\ldots e$-$34)} \\
$\ln 10$       & \texttt{(2.30258509\ldots e+00, $-$2.17075622\ldots e$-$16, $-$9.98426245\ldots e$-$33)} \\
$\varepsilon$  & $2^{-157} \approx 5.474 \times 10^{-48}$ \\
min\_normalized & $2^{-1022+2\cdot 53} \approx 1.805 \times 10^{-276}$ \\
\bottomrule
\end{tabular}
\end{center}

The minimum normalized value $2^{-916}$ ensures that all three
components remain in the normal double range.  The trigonometric
table constants $\sin(k\pi/16)$ and $\cos(k\pi/16)$ for
$k = 1, 2, 3, 4$ are similarly precomputed.


\section{C and Fortran Bindings}
\label{sec:bindings}

The QD library provides language bindings in three layers: a C API
for direct use from C code, a Fortran wrapper layer that calls the
C API via name-mangled external functions, and a Fortran~90 module
that provides operator overloading and type-safe interfaces.  The
TD extension replicates this three-layer architecture exactly.

\subsection{Three-layer architecture}

The binding architecture is:
\[
  \underbrace{\texttt{tdmodule}}_{\text{Fortran 90 module}}
  \;\xrightarrow{\text{calls}}\;
  \underbrace{\texttt{tdext}}_{\text{interface block}}
  \;\xrightarrow{\texttt{FC\_FUNC\_}}\;
  \underbrace{\texttt{f\_td.cpp}}_{\text{C++ wrappers}}
  \;\xrightarrow{\text{uses}}\;
  \underbrace{\texttt{td\_real}}_{\text{C++ class}}
\]
This mirrors the existing DD and QD paths:
\texttt{ddmodule}$\to$\texttt{ddext}$\to$\texttt{f\_dd.cpp}$\to$\texttt{dd\_real}
and
\texttt{qdmodule}$\to$\texttt{qdext}$\to$\texttt{f\_qd.cpp}$\to$\texttt{qd\_real}.

\subsection{C API (\texttt{c\_td.h} / \texttt{c\_td.cpp})}

The C wrapper represents a triple-double as a \texttt{double[3]}
array and provides functions such as:
\begin{lstlisting}
void c_td_add(const double *a, const double *b, double *c);
void c_td_mul_td_d(const double *a, double b, double *c);
void c_td_sqrt(const double *a, double *b);
void c_td_exp(const double *a, double *b);
\end{lstlisting}
Mixed-precision operations with \texttt{dd\_real} and
\texttt{qd\_real} are also provided (e.g.,
\texttt{c\_td\_add\_dd\_td}, \texttt{c\_td\_div\_td\_qd}).  The QD
interoperability functions use \texttt{to\_td\_real()} and
\texttt{to\_qd\_conversion()} for type conversion.
Self-modifying variants (\texttt{c\_td\_selfadd},
\texttt{c\_td\_selfmul\_dd}, etc.), copy functions
(\texttt{c\_td\_copy\_dd}, \texttt{c\_td\_copy\_qd}), and
comparison functions (\texttt{c\_td\_comp\_td\_qd},
\texttt{c\_td\_comp\_qd\_td}) complete the C interface.

\subsection{Fortran external interface (\texttt{tdext.f})}

The \texttt{tdext} module contains \texttt{INTERFACE} blocks
declaring the external C++ wrapper functions from
\texttt{f\_td.cpp}.  Each subroutine is declared with explicit
\texttt{INTENT} attributes and \texttt{REAL*8} arrays of length~3
for triple-double arguments.  The \texttt{PURE} attribute is
applied to all arithmetic and comparison wrappers:
\begin{lstlisting}[language=Fortran]
module tdext
  implicit none
  interface
    pure subroutine f_td_add(a, b, c)
      real*8, intent(in) :: a(3), b(3)
      real*8, intent(out) :: c(3)
    end subroutine
    ! ... 40+ additional declarations
  end interface
end
\end{lstlisting}
The \texttt{PURE} annotations enable the compiler to use these
wrappers inside \texttt{FORALL} and \texttt{DO CONCURRENT}
constructs, and are safe because the underlying C++ functions have
no side effects for non-exceptional inputs.

\subsection{Fortran wrapper (\texttt{f\_td.cpp})}

The C++ side of the Fortran interface uses \texttt{FC\_FUNC\_}
macros for portable Fortran name mangling (handling the trailing
underscore convention):
\begin{lstlisting}
#define f_td_add  FC_FUNC_(f_td_add, F_TD_ADD)
\end{lstlisting}
Each function constructs temporary \texttt{td\_real} objects from
the \texttt{double*} arguments, performs the operation, and writes
the result back via a \texttt{TO\_DOUBLE\_PTR} macro:
\begin{lstlisting}
#define TO_DOUBLE_PTR(a, ptr) \
  ptr[0] = (a)[0]; ptr[1] = (a)[1]; ptr[2] = (a)[2];

void f_td_add(const double *a, const double *b, double *c) {
  TO_DOUBLE_PTR(td_real(a) + td_real(b), c);
}
\end{lstlisting}

The Fortran wrapper also provides functions not available in the C
API, including local implementations of \texttt{floor},
\texttt{ceil}, and \texttt{aint} (truncation toward zero) for
triple-double arguments, as well as a \texttt{tdrand\_local}
function that generates random triple-double values in $[0, 1)$
by accumulating six scaled \texttt{std::rand()} calls.

\subsection{Fortran 90 module (\texttt{tdmodule})}

The \texttt{tdmodule} is the primary user-facing Fortran interface.
It \texttt{USE}s both \texttt{ddmodule} and \texttt{qdmodule}
(for the \texttt{dd\_real} and \texttt{qd\_real} types) and
\texttt{tdext} (for the external wrapper calls), then defines the
derived types and operator overloads.

\subsubsection{Derived types}

\begin{lstlisting}[language=Fortran]
type td_real
  sequence
  real*8 :: re(3)
end type td_real

type td_complex
  sequence
  real*8 :: cmp(6)     ! (1:3) = real, (4:6) = imaginary
end type td_complex
\end{lstlisting}
The \texttt{SEQUENCE} attribute ensures binary compatibility with
the C \texttt{double[3]} layout, enabling direct pass-by-reference
to the C++ wrappers without marshalling.

\subsubsection{Module constants}

\begin{lstlisting}[language=Fortran]
real*8 d_td_eps
parameter (d_td_eps = 5.47382212626881668d-48)
type (td_real) td_one, td_zero, td_eps, td_huge, td_tiny
parameter (td_one  = td_real((/1.0d0, 0.0d0, 0.0d0/)))
parameter (td_zero = td_real((/0.0d0, 0.0d0, 0.0d0/)))
parameter (td_eps  = td_real((/d_td_eps, 0.0d0, 0.0d0/)))
\end{lstlisting}
These constants mirror those in the C++ class
(\texttt{td\_real::\_eps}, etc.) and provide Fortran-idiomatic
access without calling functions.

\subsubsection{Operator overloading}

The module overloads all standard Fortran operators through
\texttt{INTERFACE} blocks.  The complete set is:

\begin{center}
\begin{tabular}{lp{7.5cm}}
\toprule
Category & Overloaded interfaces \\
\midrule
Assignment & \texttt{=} for TD$\leftrightarrow$\{double, integer,
  DD, QD, string, TD\_complex\} (21 variants) \\
Arithmetic & \texttt{+}, \texttt{-}, \texttt{*}, \texttt{/} for
  TD$\otimes$\{TD, double, integer\} and TD\_complex variants;
  \texttt{**} for power \\
Comparison & \texttt{==}, \texttt{/=}, \texttt{>}, \texttt{<},
  \texttt{>=}, \texttt{<=} for TD$\otimes$\{TD, double, integer\}
  and TD\_complex \\
\bottomrule
\end{tabular}
\end{center}

\subsubsection{Mixed-type conversions in Fortran}

The Fortran module provides explicit conversion functions via
generic interfaces:
\begin{lstlisting}[language=Fortran]
interface tdreal   ! construct td_real from other types
  module procedure to_td_d, to_td_dd, to_td_qd, to_td_i, ...
end interface
interface ddreal   ! td_real -> dd_real
  module procedure to_dd_td
end interface
interface qdreal   ! td_real -> qd_real
  module procedure to_qd_td
end interface
\end{lstlisting}
The mixed-type assignment procedures implement the same promotion
semantics as the C++ operators, but with Fortran array-slice
syntax:
\begin{lstlisting}[language=Fortran]
elemental subroutine assign_td_dd(td, dd)
  type (td_real), intent(inout) :: td
  type (dd_real), intent(in) :: dd
  td%re(1:2) = dd%re      ! copy DD components
  td%re(3) = 0.d0         ! zero-extend
end subroutine

elemental subroutine assign_qd_td(qd, td)
  type (qd_real), intent(inout) :: qd
  type (td_real), intent(in) :: td
  qd%re(1:3) = td%re      ! copy TD components
  qd%re(4) = 0.d0         ! zero-extend
end subroutine
\end{lstlisting}

\begin{remark}[Fortran vs.\ C++ conversion semantics]
The Fortran QD$\to$TD assignment (\texttt{assign\_td\_qd}) simply
copies the first three components: \texttt{td\%re = qd\%re(1:3)}.
This differs from the C++ \texttt{from\_qd\_truncate}, which
applies a 4-component renormalization before truncation.  The
Fortran version is therefore slightly less safe: if the QD value
has non-canonical component ordering (e.g., after accumulated
rounding), the truncation may lose more than
$\varepsilon_{\mathrm{td}}$ relative accuracy.  In practice this
rarely matters because QD values produced by the library maintain
the non-overlapping invariant.
\end{remark}

\subsubsection{Intrinsic function overloading}

The module overloads standard Fortran intrinsics so that
triple-double values can be used transparently:

\begin{center}
\begin{tabular}{ll}
\toprule
Intrinsic & TD overload \\
\midrule
\texttt{sin}, \texttt{cos}, \texttt{tan} & via \texttt{tdsin},
  \texttt{tdcos}, \texttt{tdtan} \\
\texttt{exp}, \texttt{log}, \texttt{log10} & via \texttt{tdexp},
  \texttt{tdlog}, \texttt{tdlog10} \\
\texttt{sqrt}, \texttt{abs} & via \texttt{tdsqrt},
  \texttt{tdabs} \\
\texttt{huge}, \texttt{tiny}, \texttt{epsilon} & return TD
  constants \\
\texttt{digits}, \texttt{precision}, \texttt{range} & return
  157, 47, 276 respectively \\
\texttt{nint}, \texttt{aint}, \texttt{floor}, \texttt{ceil}
  & rounding functions \\
\texttt{sign}, \texttt{mod}, \texttt{min}, \texttt{max}
  & utility functions \\
\texttt{random\_number} & TD-precision uniform $[0,1)$ \\
\bottomrule
\end{tabular}
\end{center}

The \texttt{td\_complex} type also overloads \texttt{conjg},
\texttt{aimag}, \texttt{exp}, \texttt{log}, and \texttt{abs}.

\subsubsection{Complex arithmetic}

The \texttt{td\_complex} type supports the same operator set as
\texttt{td\_real}, with mixed operations against \texttt{td\_real}
and \texttt{double}.  Complex arithmetic is implemented in
Fortran using the standard formulas (e.g., $(a+bi)(c+di) =
(ac-bd) + (ad+bc)i$), where each real multiplication and addition
calls the underlying TD wrappers.  Conversions to/from
\texttt{dd\_complex} and \texttt{qd\_complex} are provided.

\subsection{Quadrature test programs}

The QD library includes a suite of high-precision numerical
quadrature test programs by Bailey, implementing Gaussian
(\texttt{tquadgsq}), tanh-sinh (\texttt{tquadtsq}), and error
function (\texttt{tquaderq}) methods, including 2D variants
(\texttt{tquadgsq2d}, \texttt{tquadtsq2d}).  These programs
\texttt{USE qdmodule} and exercise the full operator overloading
infrastructure.  They serve as integration tests validating that
the Fortran module, external interface, and C++ wrappers function
correctly together across complex numerical algorithms.
Adapting these test programs from \texttt{qdmodule} to
\texttt{tdmodule} (by changing \texttt{qd\_real} to
\texttt{td\_real} and adjusting tolerance parameters from
$10^{-64}$ to $10^{-47}$) is straightforward due to the
interface-compatible module design.

\subsection{Binding layer comparison}

Table~\ref{tab:binding-layers} summarizes the parallel structure
of the DD, TD, and QD binding stacks, showing how TD fills the
middle slot with consistent naming and layout conventions.

\begin{table}[h]
\centering
\caption{Parallel structure of DD, TD, and QD binding layers.}
\label{tab:binding-layers}
\begin{tabular}{lccc}
\toprule
Layer & DD & TD & QD \\
\midrule
C++ header       & \texttt{dd\_real.h}   & \texttt{td\_real.h}   & \texttt{qd\_real.h}   \\
Inline header    & \texttt{dd\_inline.h} & \texttt{td\_inline.h} & \texttt{qd\_inline.h} \\
C wrapper        & \texttt{c\_dd.h/.cpp} & \texttt{c\_td.h/.cpp} & \texttt{c\_qd.h/.cpp} \\
Fortran wrapper  & \texttt{f\_dd.cpp}    & \texttt{f\_td.cpp}    & \texttt{f\_qd.cpp}    \\
External module  & \texttt{ddext.f}      & \texttt{tdext.f}      & (qdext)               \\
High-level module& \texttt{ddmodule}     & \texttt{tdmodule}     & \texttt{qdmodule}     \\
\midrule
Array size       & \texttt{re(2)}        & \texttt{re(3)}        & \texttt{re(4)}        \\
Complex size     & \texttt{cmp(4)}       & \texttt{cmp(6)}       & \texttt{cmp(8)}       \\
Module deps      & \texttt{ddext}        & \texttt{ddmod,qdmod,tdext} & (qdext)          \\
\texttt{PURE} wrappers & Yes             & Yes                   & Yes                   \\
\bottomrule
\end{tabular}
\end{table}

The key structural difference is that \texttt{tdmodule} depends
on both \texttt{ddmodule} and \texttt{qdmodule} (via
\texttt{USE \ldots, ONLY:}) to import the existing types for
cross-type assignment.  Neither \texttt{ddmodule} nor
\texttt{qdmodule} depends on \texttt{tdmodule}, maintaining a
strict one-directional dependency that avoids circular module
references.

\subsection{Unchanged components: \texttt{f\_dd.cpp} and
\texttt{f\_qd.cpp}}

A notable design choice is that the existing DD and QD Fortran
wrappers (\texttt{f\_dd.cpp}, \texttt{f\_qd.cpp}) are
\emph{entirely unmodified} by the TD extension.  This contrasts
with the C++ and C layers, where \texttt{dd\_real.h},
\texttt{qd\_real.h}, \texttt{c\_dd.h}, and \texttt{c\_qd.h} all
gained TD-aware operator declarations and wrapper functions
(Section~\ref{sec:mixed}).

In Fortran, all cross-type DD$\leftrightarrow$TD and
QD$\leftrightarrow$TD operations are handled entirely within
\texttt{tdmod.f} using Fortran-level type conversion followed by
same-type arithmetic.  For example, adding a \texttt{dd\_real} to
a \texttt{td\_real} in Fortran proceeds as:
(1)~the \texttt{dd\_real} is promoted to \texttt{td\_real} via
\texttt{assign\_td\_dd} (array-slice copy + zero-extension);
(2)~the TD$+$TD addition is performed by calling
\texttt{f\_td\_add} through the \texttt{tdext} interface.
This approach keeps the existing DD and QD Fortran stacks
completely stable, at the cost of an extra copy for the type
promotion.

\subsection{Complex arithmetic details}

The \texttt{td\_complex} type stores 6 doubles in a
\texttt{SEQUENCE} struct: components 1--3 for the real part and
4--6 for the imaginary part.  All complex operations are
implemented in pure Fortran within \texttt{tdmod.f} using the
underlying \texttt{td\_real} operations.  For example:

\begin{lstlisting}[language=Fortran]
type(td_real) function mul_tdc(a, b) result(c)
  ! (a1+a2*i)*(b1+b2*i) = (a1*b1-a2*b2) + (a1*b2+a2*b1)*i
  type(td_complex), intent(in) :: a, b
  type(td_real) :: a1, a2, b1, b2
  a1 = td_real(a%cmp(1:3)); a2 = td_real(a%cmp(4:6))
  b1 = td_real(b%cmp(1:3)); b2 = td_real(b%cmp(4:6))
  c%cmp(1:3) = (a1*b1 - a2*b2)%re  ! real part
  c%cmp(4:6) = (a1*b2 + a2*b1)%re  ! imaginary part
end function
\end{lstlisting}
(Simplified for presentation; the actual code uses \texttt{f\_td\_mul}
and \texttt{f\_td\_sub} calls.)  Cross-type conversions with
\texttt{dd\_complex} and \texttt{qd\_complex} are provided via
array-slice assignment (analogous to the real case), with
appropriate zero-extension or truncation.

The complex type overloads \texttt{conjg}, \texttt{aimag},
\texttt{abs}, \texttt{exp}, and \texttt{log}, providing a
complete complex arithmetic facility at triple-double precision.

\subsection{Structural parity with \texttt{ddmod.f}}

The \texttt{tdmod.f} module (1959 lines) was designed to be a
structural clone of \texttt{ddmod.f} (2077 lines, by Hida and
Bailey), with all \texttt{dd\_} prefixes replaced by \texttt{td\_}
and array dimensions changed from 2 to 3.  The module follows the
same organizational pattern: type definitions, then parameter
constants, then interface blocks for assignment, arithmetic,
comparison, intrinsics, and conversion, then the \texttt{CONTAINS}
section with all procedure implementations.

This structural parity has a practical benefit: existing Fortran
codes using DD or QD arithmetic can be migrated to TD precision
by a near-mechanical substitution of module names and type
declarations, with minimal changes to the algorithmic code itself.


\section{Operation Count Analysis}
\label{sec:flops}

This section provides a precise count of the binary64
floating-point operations (additions, subtractions,
multiplications, and FMA instructions, each counted as one
operation) required for the addition and multiplication of DD, TD,
and QD values.  The analysis is based on the actual inline
implementations in the QD library source code.

\subsection{Primitive costs}

Table~\ref{tab:primitives} summarizes the cost of the error-free
transformation primitives from Section~\ref{sec:eft}.

\begin{table}[h]
\centering
\caption{Flop counts for EFT primitives.}
\label{tab:primitives}
\begin{tabular}{lcc}
\toprule
Primitive & No FMA & FMA \\
\midrule
\textsc{QuickTwoSum} & 3 & 3 \\
\textsc{TwoSum} & 6 & 6 \\
\textsc{TwoDiff} & 6 & 6 \\
\textsc{TwoProd} & 17 & 2 \\
\textsc{ThreeSum} ($= 3 \times \textsc{TwoSum}$) & 18 & 18 \\
\textsc{ThreeSum2} ($= 2 \times \textsc{TwoSum} + 1$) & 13 & 13 \\
\textsc{QuickThreeAccum} ($= 2 \times \textsc{TwoSum}$) & 12 & 12 \\
\bottomrule
\end{tabular}
\end{table}

The \textsc{TwoProd} non-FMA cost of 17 flops comes from: $p = a
\cdot b$ (1), two \textsc{Split} calls (4 each), and the error
computation (8 flops: 4 multiplications $+$ 1 subtraction $+$ 3
additions).  With FMA, this collapses to $p = a \cdot b$ (1) and
$e = \mathrm{fma}(a, b, -p)$ (1).  This $17 \to 2$ reduction
is the dominant factor in the FMA speedup of multiplication.

The renormalization costs depend on the number of input
components and on the branch path taken. Typical-path costs are:
\textsc{Renorm3} = 6 (2 \textsc{QuickTwoSum}),
\textsc{Renorm4} = 13 ($3 + 1$ \textsc{QuickTwoSum} $+$ 1 add),
\textsc{Renorm5} = 19 ($4 + 2$ \textsc{QuickTwoSum} $+$ 1 add).

\subsection{Addition}

\begin{table}[h]
\centering
\caption{Flop counts for addition. FMA does not affect addition
  because no \textsc{TwoProd} calls are involved.}
\label{tab:add-flops}
\begin{tabular}{lcl}
\toprule
Operation & Flops & Derivation \\
\midrule
DD$+$DD (sloppy, default)
  & 11
  & $1\,\textsc{TS} + 2\,\text{add} + 1\,\textsc{QTS}$ \\
DD$+$DD (IEEE)
  & 20
  & $2\,\textsc{TS} + 2\,\text{add} + 2\,\textsc{QTS}$ \\
TD$+$TD (IEEE)
  & 53--66
  & $1\,\textsc{QTS} + (3\text{--}4)\,\textsc{QTA}
    + \text{left.} + \textsc{R4}$ \\
TD$-$TD (dedicated)
  & 56
  & $3\,\textsc{TD} + 1\,\textsc{TS}
    + 1\,\textsc{3S} + 1\,\text{add} + \textsc{R4}$ \\
QD$+$QD (sloppy, default)
  & 82
  & $5\,\textsc{TS} + 1\,\textsc{3S} + 1\,\textsc{3S2}
    + 2\,\text{add} + \textsc{R5}$ \\
\bottomrule
\end{tabular}
\end{table}

FMA has no effect on addition because the operation involves only
\textsc{TwoSum}, \textsc{QuickTwoSum}, and their derivatives, none
of which use multiplication.

For TD$+$TD, the merge-based \texttt{ieee\_add} algorithm has
data-dependent cost: if every \textsc{QuickThreeAccum} call emits
a resolved word, 3 iterations suffice (53 flops); if one call
absorbs without emitting, 4 iterations are needed (up to 66 flops).
The TD$-$TD subtraction uses a deterministic component-wise
\textsc{TwoDiff} algorithm at a fixed cost of~56.

TD has no sloppy addition variant.  This is a deliberate design
choice: the sloppy DD/QD algorithms lose roughly 1~bit of precision
compared to the IEEE variants, and this relative loss is more
significant for TD's 157-bit significand.

\subsection{Multiplication}

\begin{table}[h]
\centering
\caption{Flop counts for multiplication.}
\label{tab:mul-flops}
\begin{tabular}{lccl}
\toprule
Operation & No FMA & FMA & Derivation \\
\midrule
DD$\times$DD
  & 24 & 9
  & $1\,\textsc{TP} + 2\,\text{mul} + 2\,\text{add}
    + 1\,\textsc{QTS}$ \\
TD$\times$TD
  & 203 & 113
  & $6\,\textsc{TP} + 3\,\textsc{3S} + 4\,\textsc{TS}
    + 10\,\text{flops} + \textsc{R5}$ \\
QD$\times$QD (sloppy)
  & 207 & 117
  & $6\,\textsc{TP} + 3\,\textsc{3S} + 4\,\textsc{TS}
    + 14\,\text{flops} + \textsc{R5}$ \\
QD$\times$QD (accurate)
  & 333 & 183
  & $10\,\textsc{TP} + 3\,\textsc{3S} + 13\,\textsc{TS}
    + 18\,\text{flops} + \textsc{R5}$ \\
\bottomrule
\end{tabular}
\end{table}

Several observations:

\begin{enumerate}
\item The FMA speedup for multiplication is substantial: $2.7\times$
  for DD, $1.8\times$ for TD and QD~sloppy.  The speedup factor
  decreases with precision because the \textsc{TwoProd} cost
  ($17 \to 2$) is diluted by the accumulation and renormalization
  steps, which are FMA-independent.

\item TD$\times$TD (203 flops) is remarkably close to QD$\times$QD
  sloppy (207 flops).  The difference is only 4~flops, coming from
  the higher-order terms: TD accumulates 2 cross-products ($a_1 b_2,
  a_2 b_1$) and 4 error words ($q_0, q_3, q_4, q_5$) for 7~flops,
  while QD~sloppy accumulates 4 cross-products ($a_0 b_3, a_1 b_2,
  a_2 b_1, a_3 b_0$) and the same 4 error words for 11~flops.  The
  six \textsc{TwoProd} calls and three \textsc{ThreeSum} calls are
  \emph{identical} in both algorithms.

\item The QD accurate multiplication costs $1.6\times$ more than
  the sloppy variant (333 vs.\ 207 without FMA), primarily due to
  4~additional \textsc{TwoProd} calls and 9~additional
  \textsc{TwoSum} calls for the order-3 accumulation.

\item The cost ratio TD/DD is approximately $8.5\times$ for
  multiplication (without FMA), which is consistent with the
  $\binom{3+1}{2} = 6$ cross-products needing full error tracking
  vs.\ DD's single \textsc{TwoProd}.
\end{enumerate}

\subsection{Summary}

Table~\ref{tab:summary-flops} provides a condensed comparison,
using the default algorithm variants.

\begin{table}[h]
\centering
\caption{Summary: flop counts for default add/multiply
  (DD sloppy; TD IEEE; QD sloppy).}
\label{tab:summary-flops}
\begin{tabular}{lccccc}
\toprule
      & \multicolumn{2}{c}{Addition}
      & \multicolumn{2}{c}{Multiplication}
      & Precision \\
\cmidrule(lr){2-3} \cmidrule(lr){4-5}
Type & No FMA & FMA & No FMA & FMA & (bits / digits) \\
\midrule
DD &  11 &  11 &  24 &   9 & 106 / 31 \\
TD &  53--66 &  53--66 & 203 & 113 & 157 / 47 \\
QD &  82 &  82 & 207 & 117 & 212 / 63 \\
\bottomrule
\end{tabular}
\end{table}

The table reveals that TD provides $\approx 50\%$ more precision
than DD at roughly $5\text{--}8.5\times$ the cost, while QD provides
only $\approx 35\%$ more precision than TD at essentially the same
multiplication cost (207 vs.~203).  For applications needing 40--50
digits, TD is therefore the clear choice over QD in terms of
cost-per-digit.


\section{Conclusion}
\label{sec:conclusion}

We have described \texttt{td\_real}, a triple-double floating-point
type providing $\approx 47$ decimal digits of precision ($\approx
74\%$ of QD's 63 digits) at 49--64\% of QD's per-operation cost
(Section~\ref{sec:cost}): TD addition (sloppy, default) costs
41~flops vs.\ QD's~84 (49\%), and TD multiplication costs
116~flops (FMA) or 206~flops (no FMA) vs.\ QD's 182 or~332
(64\%/62\%).  Sloppy and accurate variants of both addition
and division are provided, following the DD/QD convention of
compile-time selection via \texttt{QD\_IEEE\_ADD} and
\texttt{QD\_SLOPPY\_DIV}.
The implementation follows the proven algorithmic framework of the QD
library---error-free transformations (\textsc{TwoSum},
\textsc{TwoProd}), renormalization, and Newton-based transcendental
functions---adapted for a three-component representation.

Key design decisions include:
\begin{itemize}
\item Overflow-safe division and square root via $2^{\pm 53}$
  rescaling when $|a_0| > 2^{969}$.
\item A dedicated TD$-$TD subtraction algorithm avoiding the
  overhead of the merge-based adder.
\item Three-stage trigonometric argument reduction with precomputed
  $\sin/\cos$ tables at $\pi/16$ intervals.
\item Cubic-convergent Newton iteration for $\log$ and $\mathrm{atan2}$.
\item Consistent mixed-precision promotion rules: TD is closed under
  double/DD/TD operations, while TD--QD operations promote to QD;
  compound-assignment from QD performs the operation in QD and
  truncates back to TD via renormalization.
\item Full C and Fortran 90 interoperability via a three-layer
  binding architecture (\texttt{tdmodule}$\to$\texttt{tdext}$\to$%
  \texttt{f\_td.cpp}$\to$\texttt{td\_real}) mirroring the existing
  DD and QD paths, with \texttt{td\_complex} support and
  overloading of all standard Fortran intrinsics.
\end{itemize}

The source code is available under a BSD 2-clause license as part
of the extended QD library.

\subsection*{Optional Branch-Free Algorithms}

libQD3 also provides compile-time optional branch-free (BF) addition and
multiplication for DD, TD, and QD arithmetic, following Kouya's transcription of
Zhang and Aiken's FPAN algorithms (arXiv:2603.14926v2).  These paths are not
enabled by default.  `QD_BF_ADD' selects BF addition, `QD_BF_MUL' selects BF
multiplication, and `QD_BF' selects both.

\begin{center}
\begin{tabular}{c|c|c|c}
Type & Operation & BF algorithm & Notes \\
\hline
DD & add & DWBFAdd, Alg. 6 & about 20 flops, error about $2u^2$ \\
DD & mul & DWBFMul, Alg. 8 & same dataflow as the existing DD multiply \\
TD & add & TWBFAdd, Alg. 11 & branch-free add, slower scalar than default \\
TD & mul & TWBFMul, Alg. 12 & branch-free multiply, exactly commutative \\
QD & add & QWBFAdd, Alg. 13 & branch-free add, slower scalar than default \\
QD & mul & QDBFMul, Alg. 14 & branch-free multiply, exactly commutative
\end{tabular}
\end{center}

\begin{thebibliography}{9}
\bibitem{HLB2001}
Y.~Hida, X.~S.~Li, and D.~H.~Bailey,
``Algorithms for quad-double precision floating point arithmetic,''
in \textit{Proceedings of the 15th IEEE Symposium on Computer
Arithmetic}, 2001, pp.~155--162.

\bibitem{QDlib}
Y.~Hida, X.~S.~Li, and D.~H.~Bailey,
``Library for double-double and quad-double arithmetic,''
LBNL Technical Report, 2007.
\url{https://www.davidhbailey.com/dhbpapers/qd.pdf}

\bibitem{Dekker1971}
T.~J.~Dekker,
``A floating-point technique for extending the available precision,''
\textit{Numerische Mathematik}, vol.~18, no.~3, pp.~224--242, 1971.

\bibitem{Knuth1997}
D.~E.~Knuth,
\textit{The Art of Computer Programming}, vol.~2:
\textit{Seminumerical Algorithms}, 3rd ed.
Addison-Wesley, 1997.

\bibitem{Priest1991}
D.~M.~Priest,
``Algorithms for arbitrary precision floating point arithmetic,''
in \textit{Proceedings of the 10th IEEE Symposium on Computer
Arithmetic}, 1991, pp.~132--143.

\bibitem{Shewchuk1997}
J.~R.~Shewchuk,
``Adaptive precision floating-point arithmetic and fast robust
geometric predicates,''
\textit{Discrete \& Computational Geometry},
vol.~18, no.~3, pp.~305--363, 1997.

\bibitem{Boldo2008}
S.~Boldo and G.~Melquiond,
``Emulation of a FMA and correctly-rounded sums: proved algorithms
using rounding to odd,''
\textit{IEEE Transactions on Computers},
vol.~57, no.~4, pp.~462--471, 2008.
\end{thebibliography}

\end{document}
